\chapter[Beauty: Object or Beholder?]{Is Beauty in the Object or the Beholder?}
\section{A Stone in a Drawer}
First, pick up an object. Not a museum treasure: something from your drawer.

Many drawers hold one: a found stone. Perhaps on a primary-school outing, beside a river, mountain path, or retreating sea foam, you crouched and chose it from thousands. Gray with a white band, perhaps, or deep blue-green worn round by water. Hold it long enough and it takes your warmth. You brought it home, washed it, and put it away. For years, every clearing-out brought a promise to discard it, never kept.

Now imagine a classmate visiting, finding it, turning it over, and asking, "Isn't it just a stone? Why pick it up?"

You probably cannot explain. You might say "beautiful," then stumble at "what's beautiful about it?" You cannot name its composition, age, or value. It is indeed worthless; thousands similar remain on the riverbank. Yet you remember crouching: among all those stones, this one made you pause. That pause was real.

A geologist can provide a complete report: mainly quartz, containing mica, specified density, hardness, formation date, and duration of water erosion. Every item is objective: identical instruments yield identical results for anyone. Search the report, though, for "beautiful." It is absent. Instruments measure everything except what made you crouch.

Where does beauty live? If inside the stone like density and hardness, why cannot your classmate detect it? Is something wrong with him, or are his instruments insufficient? If merely your feeling, a private report like hunger, why did crouching feel like \textbf{discovery}, not \textbf{invention}? Your attention faced the stone, not your stomach.

A stone blocks both easy exits. Begin here: \textbf{is beauty an object's property or a beholder's feeling?}

\section{Two Camps in the Gallery}
Enter a scene with me.

Saturday afternoon, the municipal art museum. Light rain outside, air-conditioning inside, polished floors softly squeaking under trainers. Few visitors, lowered voices, as in a great library.

Gallery C holds ceramics labeled Song dynasty. Gentle lighting isolates a small celadon bowl in glass. Its indescribable blue-green leans gray, then green; a fine crack climbs from rim to base.

Gallery D holds one wall-sized painting: black and gray split diagonally by dark red, shapeless at its edges, seemingly splashed or scraped. The label gives only title and year, no explanation.

Four groups occupy the galleries.

You and your classmate A-Can. He stands motionless before the bowl for three full minutes. Urged onward, he waves without looking away. You inspect it for ten seconds and privately think: just a bowl, even cracked; the cafeteria has stacks. You feel too embarrassed to say so.

A middle-aged couple faces the huge painting. Hands on hips, the husband cannot hide anger even in a low voice. "This is painting? Leftover decorating paint thrown on a wall would do it. They built a gallery for this?" His wife does not answer. Arms folded, head tilted, she watches for a long time, then softly says, "Please be quiet. Let me look." Angrier, he asks what she sees. She thinks. "I can't say. But please be quiet."

In a corner, an elderly visitor with a magnifying glass slowly reads tiny inscriptions and seals on a landscape painting as though reading a long letter. A guard by the doorway struggles to suppress a yawn.

The strangest thing: nobody disputes facts. Nobody doubts the bowl's dynasty, misreads the red's darkness, or measures different dimensions. \textbf{Everyone agrees at the level of what is there.} Yet "is it good, beautiful, worth displaying?" instantly splits the room. One person cannot move from an object; another watches in bewilderment, as though seeing someone stare at empty air.

Both feel justified. The husband thinks "art" has intimidated his wife, an emperor's-new-clothes game where dissent means ignorance. The wife thinks he is not looking at all; he is angry at "art," not attending to the red.

And you? Your puzzlement at celadon and his anger at the painting are the same phenomenon. You share a camp, merely in different galleries.

\section[Does Personal Taste End Discussion?]{Does "Everyone Has Their Taste" End Discussion?}
Lay out the conflict before answering.

We have a ready truce: "Different tastes for different people." Taste is private, like preferring carrots or greens; nobody can interfere, discussion ends. It sounds generous, mature, peaceable. Take it seriously: it advances a complete position, \textbf{beauty is subjective, so arguing about taste wastes time}.

Do not sign yet. Parts of this truce fail to fit.

First, aesthetic arguments do not resemble taste reports. "I don't eat coriander." "Then more spring onion." Finished, without passion. The gallery husband is angry, the wife stubborn; each thinks the other \textbf{wrong}, not merely \textbf{different}. If aesthetic judgment were a private report like hunger or ticklishness, "you're wrong" would not apply. Nobody normally says, "You say you're hungry? Wrong, you're not." Yet beauty continually elicits "wrong," "you missed it," "bad taste." Our behavior betrays us: verbally we sign the truce while treating the issue like a contested fact.

Second, if personal taste ends matters, "learning" should be expelled from aesthetics. A-Can was not born able to linger over celadon. He once found museums as boring as you. A teacher later had him copy ceramics, handle clay, and learn how the kiln produced that crack and how many firings potters waited for such patterns. His eyes changed. \textbf{Seeing can be trained}: aesthetics is not entirely factory settings. Conversely, if beauty is wholly trained, a teacher's implanted answer, how does learning differ from memorization? Are those three minutes yours or the teacher's?

Third, most stubbornly, some things produce remarkable agreement. When sunset blazes, people on grasslands, islands, and city streets, of every skin color, ancient and modern, generally pause. If beauty is private taste, whence such widespread agreement? Carrots and greens have never united humanity like this.

So a deeper question emerges than who is right:

\textbf{What are we doing when we say, "How beautiful"?} Reporting private pleasure beyond others' judgment? Identifying a property anyone can check? Or doing a third thing we cannot quite describe?

This is not idle after-dinner conversation. It affects whether schools teach art and what they teach; who can call an old neighborhood ugly enough to demolish; whether a tenfold clothing price reflects fabric or appearance; and whether your classmate's dismissal of the stone deserves a shrug or an invitation to look until something appears.

Before continuing, take a position in the gallery. Which spouse seems nearer the truth, or is the question mistaken?

\section{Both Sides Hold Strong Cards}
Give both sides their strongest reasoning. This exercise is "steelmanning": strengthen the opponent rather than constructing a straw figure easily knocked down. An unfair victory is no victory.

\textbf{First position: beauty belongs to things themselves.}

Play this side's full hand. Its foundational cards are ancient and firm.

Tradition says Pythagoras, passing a smithy about twenty-five centuries ago, heard some hammer combinations pleasing, others harsh. Studying strings afterward, he found half a string's length produced an octave, harmonizing closely with the original, while two-thirds also yielded consonance. Thus \textbf{simple whole-number ratios lay behind agreeable sound}. Greeks were astonished: harmony was a numerical relationship in string lengths, unchanged by who measured, not private taste. Beauty's secret seemed proportion, symmetry, order: balanced faces, proportioned buildings, integer-ratio intervals. This shaped European art and architecture for over two millennia.

Modern cards follow. Psychological research commonly observes relationships between facial attractiveness and symmetry, with unexpectedly high agreement across cultures. Sunset belongs here too: purely private preference struggles to explain shared response. This camp asks whether every agreement is coincidence. Judgments are plainly \textbf{triggered} by properties such as symmetry, proportion, and harmony, which belong to objects rather than minds.

\textbf{Second position: beauty belongs to the beholder's feelings.}

Its cards are equally strong. Eighteenth-century Scottish philosopher David Hume discussed beauty as residing not in things but in observing minds. One object can appear entirely different to different people without anyone misperceiving; beauty lives in feeling.

Evidence? Cultural variation. Japanese wabi-sabi prizes age, imperfection, asymmetry: a worn, dark-glazed, chipped tea bowl may surpass a flawless new one. Kintsugi repairs ceramics with gold-bearing lacquer, highlighting rather than concealing fractures as brilliant lines; breakage's history becomes beauty. Eighteenth-century European courts seeking immaculate perfection might find this baffling. Song literati might find Rococo halls crowded with gilding, scrollwork, and mirrors garishly vulgar. \textbf{Change the culture and the same object's evaluation reverses.} A property like density should not reverse that way.

Sharper examples follow. Centuries ago, elites in Japan and parts of Southeast Asia prized blackened teeth as attractive signs of status. Other cultures and periods preferred pale skin, fuller bodies, or thinness, turning direction within decades. Object-determined taste cannot explain such collective reversals. Beauty situated in beholders, periods, classes, and fashion can.

\textbf{Do not choose between them yet.}

Each holds cards beyond the other's reach. Object-centered accounts explain agreement but struggle with cultural difference; beholder-centered accounts do the reverse. A third route is \textbf{relational understanding}: beauty lies neither in the object's pocket nor wholly in the mind, but in their encounter. Properties, symmetry, cracks, dark red, meet \textbf{this} viewer's body, experiences, training, and culture. The stone has an attention-catching texture, but you must be the person crouching beside the river. Comprehensive though this sounds, it has a problem: if it explains everything, how can it decide which side wins an argument? We will return to that.

\section{A Bridge for Thought: Four Statements Hidden in "Beautiful"}
Can a portable tool replace volume and seniority in aesthetic disputes? Yes. When hearing or saying "beautiful" or "what's beautiful about that?", separate four statements and locate the speaker.

\textbf{First: a bodily report.} "Seeing it moved me." This is an internal fact like hunger or ticklishness, \textbf{neither mistaken nor needing to persuade}. The wife's body responds before the painting; nobody may deny that response as wrong.

\textbf{Second: personal preference.} "I like it." Preference adds history and habit to bodily response: repeated spicy food makes spice bearable; familiar music sounds agreeable. Sweet versus salty need not be debated. A-Can likes celadon; you like trainers' lines. That is statement two.

\textbf{Third: judgment about an object.} "It is beautiful." Grammar changes: "I" becomes "it." You judge \textbf{that thing}, no longer merely report yourself, inviting or even demanding that others see it too. Gallery disputes begin in the slide from two to three. The husband's anger concerns not his wife's preference but her supposed mistaking rubbish for treasure. Her "be quiet" implies "look again; you will see."

\textbf{Fourth: reasons for judgment.} "It is beautiful because..." Discussion properly begins here, since reasons can be right or wrong. "Because it is genuine Song ware" can be overturned archaeologically. "Because the crack lets the bowl's color breathe" points toward the object. Others can follow your finger and say either "I still don't see it" or "Now that you say it, I do."

This clears confused accounts. Disputes last all night because speakers inhabit different statements: one reports being moved, another denies objective beauty. Different subjects masquerade as disagreement. When a parent asks what is good about a song, distinguish denial of your response, which needs no approval, from argument over the object, which can advance to shared reasons.

The tool also maps "everyone has their taste." It fully governs bodily reports and preferences, but not judgments and reasons. Once someone judges an object and offers reasons, discussion begins and the truce lapses. The saying is not wrong, only misplaced: a ceasefire covering half the territory.

\section{Zhuangzi's Fishpond}
Read a short primary passage. Twenty-three centuries ago, a Chinese thinker pushed beauty's location to an extreme.

The Zhuangzi's "Discussion on Making All Things Equal" says:

\begin{quote}
Mao Qiang and Lady Li are beautiful to people. Fish seeing them dive deep, birds fly high, and deer run away. Which of these four knows the world's true beauty?
\end{quote}
Mao Qiang and Lady Li were legendary beauties. Humans admired them, while fish submerged, birds fled upward, and deer bolted. Which observer knew beauty's proper standard?

Locate this on our map. It challenges extreme object-centered beauty: if beauty were an intrinsic property like height, every seeing observer should register the same result. Zhuangzi offers one object and four observers with four responses. Fish and birds may not even find the women ugly. A fish's experiential world may have no "beautiful woman" category; an admired human figure is simply an approaching large object, a danger. \textbf{Beauty is not a signal emitted by an object alone, but a connection between object and observer; without a connection, signal becomes noise.}

Zhuangzi goes deeper than "different tastes," which contrasts human preferences. He asks who sets \textbf{the standard itself} for legitimate beauty. Humans admire Mao Qiang and proclaim universal beauty, yet this universe seats no fish. He dismantles not a particular judgment but the quiet promotion of human feeling into cosmic law. Every gallery "what's beautiful?" and "you don't understand" similarly declares its own connection the only legitimate one.

Keep one essential reservation. Giving fish, birds, and deer their own experience acknowledges the observer's reality and weight. The passage does not prove anything goes because beauty is purely subjective; it establishes that \textbf{beauty cannot exist apart from a beholder}. Objects must meet someone capable of feeling. This leads toward relational understanding. Zhuangzi does not consider all encounters equivalent either: the entire book trains another way of seeing. Its stance opens perception rather than gives up.

The perspective is remarkably modern. Biology finds enormous variation in eyes and sensitivity: bees see ultraviolet floral patterns humans miss; many birds inhabit richer color worlds. Assuming human sight shows the world as it really is fails scientifically too. Zhuangzi knew no ultraviolet, yet a fishpond led him to the same insight: an observer's position determines the world received. An ancient fable and modern science share a starting line because the question has deep roots.

\section{One Sunset, Three Explanations}
Return to agreement: why do almost all people pause at a blazing sunset? We can now compare genuine alternatives. Separate evidence, substantial cross-cultural agreement on sunsets, consonance, and certain faces; explanation, where theories compete; contemporaries' accounts, perhaps gods painting heaven; and value judgments. We compare explanation.

\textbf{First: beauty is inherited preference, a biological account.}

Bodies shaped by millions of years retain survival-manual preferences. Open landscapes with water and clear views feel good because ancestors preferring them more easily found food and detected danger, passing preferences onward as beauty. Symmetrical faces may signal healthy development. Sunsets may promise fine weather or recall fires at the center of human nights, embedding warm colors in security. This addresses global agreement with testable mechanisms. Its limits: it explains comfort and visual ease better than learned wabi-sabi, kintsugi, or disturbing tragedy. More fundamentally, innate preference cannot establish what \textbf{should} count as beautiful. Ancestors favored sugar and fat too; cream cake does not therefore deserve a nutrition prize. Innate preference is a defendant to question, not a judge.

\textbf{Second: culture teaches beauty.}

Shift from bodies to cultivation. Celadon pleases because your culture enthrones Song ceramics; textbooks, documentaries, and museum labels calibrate your eyes. Wabi-sabi's golden cracks or an era of blackened teeth would calibrate them differently. Fashion supplies quick evidence: flared trousers despised twenty years ago become stylish again for the same person. The object remains; learned judgment reverses. This account takes trainable seeing seriously and explains aesthetic links to class and identity: what you appreciate signals who you are. But purely learned accounts struggle with shared responses, including experiments finding infants prefer consonant to dissonant sounds. Calling every agreement coincidence or learning is strained. And who taught culture? Did its original standards fall from the sky?

\textbf{Third: beauty emerges from three things meeting.}

Take part of each account. Objective properties, symmetry, proportion, color, rhythm, cracks, exist independently of wishes. Yet different bodies, experiences, and training respond differently. Beauty belongs to the \textbf{relationship}. Biology supplies factory settings; experience engraves personal grooves, such as your third-grade outing in the drawer's stone; culture overlays a shared map. Together they form you before the sunset. This accommodates both agreement through shared bodies and Earth, and difference through experience and culture. Its limitation is honesty about its role: comprehensive as \textbf{description}, nearly silent as \textbf{judge}. It can tell the arguing couple their relationships differ, not who is right. Many most want precisely that: can one aesthetic judgment have better reasons? Next we face it.

Each explanation grasps part, not all. Their comparison seeks no champion but reveals three different origins beneath a casual "beautiful."

\section{A Deeper Challenge: Pleasure without Use}
Now introduce the weightiest theory. Late-eighteenth-century philosopher Immanuel Kant dismantled and rebuilt beauty's location in the Critique of Judgment. The original is difficult; we use only a helpful framework, entirely paraphrased.

Kant begins with an overlooked peculiarity. "This flower is beautiful" arises from feeling, not a ruler's measurement, yet sounds like a claim everyone should accept. Rather than limiting it to ourselves, we expect, even \textbf{demand}, agreement. Aesthetic judgment is subjective, unsupported by concepts or proof, yet speaks with an "ought to be universal" tone. How can one sentence combine these?

His first distinction is "disinterestedness." An estate agent sees an old house's sale price; a carpenter its beams' strength; a visitor simply its beauty. The first two look through the house toward money and use; the third stops at the object. Only pleasure without instrumental purpose is aesthetic pleasure. This explains why a painting's hundred-million price need not move its true viewer and why your classmate's "worthless stone" cannot hurt you: different measures apply. Usefulness and price are invalid currencies here.

The second step is more ambitious. Although no formula proves a flower beautiful, demanding agreement presupposes broadly shared human capacities. Kant calls this "common sense": we assume similar perceptual equipment, so explaining and pointing out our way of seeing might let another see too. The theory accommodates both camps. Judgment points to an object's properties, but without a sensitive beholder those properties do not become beauty. Argument need not be wasted because common sense makes reasons communicable. The wife's "let me look" is not defeated inarticulacy but an unfinished invitation to share a way of seeing.

Kant holds another card, supplying an easily missed point in the brief: \textbf{beauty is not prettiness, and art need not bring comfort}.

He distinguishes the beautiful from the sublime. The beautiful is small, gentle, relaxing: sunset, porcelain, rounded stone. The sublime comes from immensity and terror: a storm crossing the sea viewed from a cliff, the Milky Way overhead. Smallness and fear give way to strange exhilaration: small though I am, my mind can contain this vastness. Comfortable? Not at all. Yet few deny the aesthetic experience, perhaps a deeper one. Earlier still, Aristotle's Poetics described tragedy arousing pity and fear to bring their catharsis. Audiences leave devastating drama in tears yet feel cleansed. Humanity has needed uncomfortable tragedy for two millennia. Picasso's Guernica depicts a bombed town through broken, twisted, crying black-and-white forms, neither pretty nor comfortable, yet among war's most powerful indictments.

Together these dismantle a simplistic equation: \textbf{beauty \ensuremath{\ne} pleasure \ensuremath{\ne} comfort \ensuremath{\ne} prettiness}. Prettiness relaxes, but beauty's territory includes awe, heartbreak, and unease. Art does not merely serve the world sugar. Discomfort does not prove failure; comfort alone need not exhaust value. When someone dismisses an uncomfortable painting as non-art, distinguish their valid bodily report from an invalid inference. Between discomfort and denied beauty stands Kant's whole distinction between beautiful and sublime.

Admit the limits. Kant's shared equipment meets Zhuangzi's reminder of enormous cultural variation. Someone unfamiliar with wabi-sabi has no preinstalled channel for admiring cracks. Common sense may exist, with coverage smaller than Kant assumed and greater than radical private taste allows. Its boundary remains open.

\section{Who Finds Beauty for You in the Filter Age?}
This ancient question now lives more actively than ever in your phone.

Beauty filters smooth skin, enlarge eyes, and narrow faces at a tap. They disguise \textbf{a particular training of the beholder}, certain symmetry, color, proportions, as \textbf{the object's own nature}, as though a face ought naturally to look that way. They smuggle the third statement, this face is beautiful, into the first bodily response, deciding before you judge. A loop follows: standardized faces retrain your eyes until only that template looks right. "Everyone has their taste" becomes "everyone likes what the algorithm teaches." Difference has not vanished but been flattened.

Then consider songs growing pleasing through repetition. Psychology repeatedly finds mere exposure often increases liking. A noisy first impression becomes a tune hummed after dozens of short-video repetitions. This is ordinary relational change, but platform-controlled exposure makes preference's origins worth tracing. Was it your spontaneous riverbank pause, or a stone placed before you repeatedly until you crouched? In Kantian terms, the danger is not mistaken judgment but \textbf{mistaking delivery for judgment}.

A more serious echo follows. Aesthetic standards have bodily weight. Chinese foot-binding persisted nearly a millennium under beauty's name, broken feet becoming marriage-market credentials. European corsets in roughly corresponding periods displaced organs, also for beauty. These women were not simply masochists; many sincerely admired the result. That is precisely the warning: \textbf{a repeatedly calibrated cultural standard can become the body's own desire, defended even by its victims}. Understanding sincere admiration does not validate the standard. Foot-binding is gone today, yet fixation on weighing under fifty kilograms, minors seeking leg liposuction, and replicated standard faces continue the logic. Aesthetic discussion is not harmless small talk; here it touches the ground.

Do not paint the present entirely gray. A small-town teenager can freely view Dunhuang murals, Louvre collections, and kintsugi hands, with a cultural-learning radius beyond any historical emperor's. Algorithms flatten difference while open collections nourish it. "Tools are neutral" is lazy, but direction remains contestable. Begin by tracing each experience of being moved.

\section{Your Question: Can a Vote Decide Beauty?}
Return to the drawer's stone.

Your classmate calls it "just a stone." You now have a map: which statement is his, which can you answer? You need not defend the first feeling. Or bring him to the desk lamp and show the white band curling around it, offering fourth-statement reasons and an invitation. He may still see just a stone. This time you have not merely failed to agree; you have genuinely discussed it.

Both camps have their strongest cards: agreement, proportion, bodies on one side; culture, historical reversals, and Zhuangzi's fishpond on the other. Relational understanding takes both in but cannot declare a winner. Kant offers common sense whose reach nobody measures precisely.

So ask: \textbf{can aesthetic judgments truly be ranked?}

Suppose your class votes on a picture for its wall. Most choose a finely printed celebrity poster; three choose a gray print whose appeal they cannot explain. Majority rule governs decoration. But can votes determine not just display, but \textbf{beauty}? If the three say, "You haven't really looked yet," what authorizes them? If right, beauty resists voting; if wrong, it retreats to private taste. Who then explains our agreements on sunsets, harmonious sounds, and the moment almost everyone pauses?

Give reasons for your answer. Weigh agreement, cultural difference, trained seeing, and the intransferable weight of your own riverbank pause.

There is no standard answer. But from today, whenever you say "beautiful" or "what's beautiful about that?", listen for the layer beneath your words.
